\section{闭集}

记号约定:$X$是拓扑空间,$A\subseteq X$.

\begin{definition}{闭集}{}
	若$X\setminus A$是$X$中的开集,则称$A$是$X$中的闭集.
\end{definition}

\begin{definition}{开闭集}{}
	$X$中既是开集又是闭集的集合称为开闭集.
\end{definition}

显然,$\emptyset$和$X$一定是开闭集.

\begin{proposition}{闭集公理}{}
	\begin{enumerate}
		\item ($\mathrm{O}_{1}^{\prime}$)$X$中任意一族闭集的交还是$X$中的闭集.
		\item ($\mathrm{O}_{2}^{\prime}$)$X$中任意两个闭集的并还是$X$中的闭集.
	\end{enumerate}
\end{proposition}

\begin{example}{$\Q$中的闭集}{}
	在有理直线$\Q$中:
	\begin{enumerate}
		\item 左开右闭区间$\left[a,+\infty\right)$和$\left(-\infty,a\right]$都是闭集($a\in\Q$),因为二者的补集分别为$\left(-\infty,a\right)$和$\left(a,+\infty\right)$,这都是开集.
		\item 有界闭区间$\left[a,b\right]$是闭集(其中$a,b\in\Q$且$a<b$),因为$\left[a,b\right]=\left(-\infty,b\right]\cap\left[a,+\infty\right)$.
		\item (在典范嵌入的意义下)$\Z$是闭集,因为$\Q\setminus\Z=\bigcup\limits_{n\in\Z}\left(n,n+1\right)$.
	\end{enumerate}
\end{example}

\begin{definition}{闭覆盖}{}
	设$\left(F_{i}\right)_{i\in I}$是$X$中一族闭集.若该闭集族是$A$的覆盖,则称之为$A$在$X$中的闭覆盖.
\end{definition}

\begin{proposition}{同胚的闭集刻画}{}
	设$f:X\to Y$是双射,则$f$是同胚当且仅当下列条件成立:
	\begin{enumerate}
		\item $X$中每个闭集在$f$下的像都是$Y$中的闭集.
		\item $Y$中每个闭集在$f^{-1}$下的像都是$X$中的闭集.
	\end{enumerate}
\end{proposition}